\documentclass[11pt]{article}

\usepackage[margin=1in]{geometry}
\usepackage{array}
\usepackage{longtable}
\usepackage{booktabs}
\usepackage{parskip}
\usepackage{enumitem}
\usepackage{url}
\usepackage{graphicx}
\usepackage{amsmath,amssymb}
\usepackage{float}

\setlength{\LTcapwidth}{\textwidth}

% Same look as \titleformat{\section}{\large\bfseries}{}{0em}{} from A1, without needing titlesec
\makeatletter
\renewcommand\section{\@startsection{section}{1}{\z@}{1.2em}{0.5em}{\large\bfseries}}
\renewcommand\subsection{\@startsection{subsection}{2}{\z@}{0.6em}{0.2em}{\normalsize\bfseries}}
\makeatother

% Column types for the tables
\newcolumntype{L}[1]{>{\raggedright\arraybackslash}p{#1}}

\graphicspath{{figures/}}
\newcommand{\code}[1]{\texttt{#1}}

\title{\vspace{-3em}{CSC490 Assignment A4: RL-ing nanochat\vspace{-1em}}}

\author{
  Prithvi Seran (1009973343) \and
  Vraj Shah (1009102695) \and
  Tyseer Toufiq (1009364064) \and
  Joshua Crisologo (1009860438)
}
\date{}

\begin{document}


\maketitle
\vspace{-2.5em}
\begin{center}
\small Code, logs and scripts: \url{https://github.com/Vraj2004/nanochat-490}
\end{center}

\section*{Part 1: GRPO and RL Review}

GRPO \cite{grpo} is a variant of PPO \cite{ppo} that replaces the learned value function with a group baseline. For each question $q$ it samples $G$ outputs from the old policy, scores each with a reward $r_i$, and maximizes
\begin{equation}
J_{\text{GRPO}}(\theta) = \frac{1}{G}\sum_{i=1}^{G}\frac{1}{|o_i|}\sum_{t=1}^{|o_i|}
\Big\{\min\!\big[\rho_{i,t}\hat{A}_i,\ \text{clip}(\rho_{i,t},1-\varepsilon,1+\varepsilon)\,\hat{A}_i\big]
-\beta\, D_{\text{KL}}\big[\pi_\theta\,\|\,\pi_{\text{ref}}\big]\Big\},
\label{eq:grpo}
\end{equation}
where $\rho_{i,t} = \pi_\theta(o_{i,t}\mid q,o_{i,<t})\,/\,\pi_{\theta_{\text{old}}}(o_{i,t}\mid q,o_{i,<t})$ and $\hat{A}_i = (r_i-\text{mean}(\mathbf{r}))/\text{std}(\mathbf{r})$.

nanochat's \code{scripts/chat\_rl.py} \cite{nanochat} keeps the group sampling and the mean baseline and removes almost everything else. Its docstring puts ``GRPO'' in quotes and calls the result ``more similar to just REINFORCE'' \cite{reinforce}. For each GSM8K \cite{gsm8k} question it samples 16 completions and maximizes
\begin{equation}
J_{\text{nanochat}}(\theta) = \frac{1}{N}\sum_{i}\sum_{t} m_{i,t}\,\big(r_i-\text{mean}(\mathbf{r})\big)\,\log\pi_\theta(o_{i,t}\mid q,o_{i,<t}),
\label{eq:nanochat}
\end{equation}
where $m_{i,t}$ is 1 only for tokens the model sampled (0 for the prompt, padding and inserted calculator outputs) and $N$ is the number of such tokens in the forward pass. Table~\ref{tab:grpo} lists the four changes.

\begin{table}[H]
\centering
\small
\renewcommand{\arraystretch}{1.15}
\caption{GRPO in DeepSeekMath \cite{grpo} against \code{chat\_rl.py} at its default settings.}\label{tab:grpo}
\begin{tabular}{L{2.7cm} L{5.6cm} L{6.5cm}}
\toprule
\textbf{Component} & \textbf{GRPO} & \textbf{nanochat} \\
\midrule
Surrogate & PPO ratio $\rho$ against $\pi_{\theta_{\text{old}}}$, clipped & $\log\pi_\theta$ times the advantage; one optimizer step per 256 fresh rollouts \\
Trust region & KL penalty to a reference policy, $\beta = 0.04$ & None; the learning rate starts at 5\% of its base value and decays linearly to zero over one epoch (467 steps) \\
Advantage & Group mean subtracted, then divided by the group std & Group mean subtracted only \\
Length normalization & Mean over each output's tokens, then mean over outputs & Sum over all sampled tokens, divided by their count \\
\bottomrule
\end{tabular}
\end{table}

\textbf{Why Karpathy may have made these changes.} nanochat is meant to be minimal, and each removed piece does little in this setting. Every batch of rollouts is used for one gradient step and then discarded, so $\rho = 1$, the clip is inactive, and the gradient of the clipped term in Eq.~\ref{eq:grpo} is $\hat{A}_i\nabla\log\pi_\theta$, which is REINFORCE with a baseline. The KL penalty costs a frozen copy of the model and an extra forward pass. DAPO \cite{dapo} and Dr.~GRPO \cite{drgrpo} drop it too, and here a small learning rate that decays to zero within one epoch already keeps the policy near the SFT model. Dividing by the std up-weights questions whose samples nearly all agree: with a 0/1 reward, a question with 1 of 16 samples correct gets about twice the weight of one with 8 of 16. Dr.~GRPO calls this a question-level difficulty bias and removes the division as well. Averaging within each output gives each token of a long answer less weight than a token of a short one, and DAPO replaces it with the token-level average that nanochat uses. One limit carries over from GRPO: a question whose 16 samples all get the same reward has zero advantage and adds nothing to the gradient. In our replication (Part 3) that was 46\% of the training questions, which the shaped rewards in Part 4 are aimed at.

% Parts 2-5 go here.

\begin{thebibliography}{99}

\bibitem{nanochat}
A. Karpathy, ``nanochat: A tiny chatbot arena and training harness,'' GitHub, 2025. \url{https://github.com/karpathy/nanochat/discussions/481}

\bibitem{grpo}
Z. Shao et al., ``DeepSeekMath: Pushing the Limits of Mathematical Reasoning in Open Language Models,'' arXiv:2402.03300, 2024.

\bibitem{ppo}
J. Schulman et al., ``Proximal Policy Optimization Algorithms,'' arXiv:1707.06347, 2017.

\bibitem{reinforce}
R. J. Williams, ``Simple Statistical Gradient-Following Algorithms for Connectionist Reinforcement Learning,'' \textit{Machine Learning}, vol.~8, 1992.

\bibitem{gsm8k}
K. Cobbe et al., ``Training Verifiers to Solve Math Word Problems,'' arXiv:2110.14168, 2021.

\bibitem{dapo}
Q. Yu et al., ``DAPO: An Open-Source LLM Reinforcement Learning System at Scale,'' arXiv:2503.14476, 2025.

\bibitem{drgrpo}
Z. Liu et al., ``Understanding R1-Zero-Like Training: A Critical Perspective,'' arXiv:2503.20783, 2025.

\end{thebibliography}

\end{document}
